% Options for packages loaded elsewhere
\PassOptionsToPackage{unicode}{hyperref}
\PassOptionsToPackage{hyphens}{url}
\PassOptionsToPackage{dvipsnames,svgnames,x11names}{xcolor}
\PassOptionsToPackage{space}{xeCJK}
\documentclass[
  chinese,
  10pt,
  a4paper,
]{article}
\usepackage{xcolor}
\usepackage[margin=2cm]{geometry}
\usepackage{amsmath,amssymb}
\setcounter{secnumdepth}{-\maxdimen} % remove section numbering
\usepackage{iftex}
\ifPDFTeX
  \usepackage[T1]{fontenc}
  \usepackage[utf8]{inputenc}
  \usepackage{textcomp} % provide euro and other symbols
\else % if luatex or xetex
  \usepackage{unicode-math} % this also loads fontspec
  \defaultfontfeatures{Scale=MatchLowercase}
  \defaultfontfeatures[\rmfamily]{Ligatures=TeX,Scale=1}
\fi
\usepackage{lmodern}
\ifPDFTeX\else
  % xetex/luatex font selection
  \setmainfont[]{Times New Roman}
  \ifXeTeX
    \usepackage{xeCJK}
    \setCJKmainfont[FontIndex=3,BoldFont={/System/Library/AssetsV2/com_apple_MobileAsset_Font8/86ba2c91f017a3749571a82f2c6d890ac7ffb2fb.asset/AssetData/PingFang.ttc},BoldFeatures={FontIndex=11}]{/System/Library/AssetsV2/com_apple_MobileAsset_Font8/86ba2c91f017a3749571a82f2c6d890ac7ffb2fb.asset/AssetData/PingFang.ttc}
  \fi
  \ifLuaTeX
    \usepackage[]{luatexja-fontspec}
    \setmainjfont[FontIndex=3,BoldFont={/System/Library/AssetsV2/com_apple_MobileAsset_Font8/86ba2c91f017a3749571a82f2c6d890ac7ffb2fb.asset/AssetData/PingFang.ttc},BoldFeatures={FontIndex=11}]{/System/Library/AssetsV2/com_apple_MobileAsset_Font8/86ba2c91f017a3749571a82f2c6d890ac7ffb2fb.asset/AssetData/PingFang.ttc}
  \fi
\fi
% Use upquote if available, for straight quotes in verbatim environments
\IfFileExists{upquote.sty}{\usepackage{upquote}}{}
\IfFileExists{microtype.sty}{% use microtype if available
  \usepackage[]{microtype}
  \UseMicrotypeSet[protrusion]{basicmath} % disable protrusion for tt fonts
}{}
\makeatletter
\@ifundefined{KOMAClassName}{% if non-KOMA class
  \IfFileExists{parskip.sty}{%
    \usepackage{parskip}
  }{% else
    \setlength{\parindent}{0pt}
    \setlength{\parskip}{6pt plus 2pt minus 1pt}}
}{% if KOMA class
  \KOMAoptions{parskip=half}}
\makeatother
\usepackage{longtable,booktabs,array}
\newcounter{none} % for unnumbered tables
\usepackage{calc} % for calculating minipage widths
% Correct order of tables after \paragraph or \subparagraph
\usepackage{etoolbox}
\makeatletter
\patchcmd\longtable{\par}{\if@noskipsec\mbox{}\fi\par}{}{}
\makeatother
% Allow footnotes in longtable head/foot
\IfFileExists{footnotehyper.sty}{\usepackage{footnotehyper}}{\usepackage{footnote}}
\makesavenoteenv{longtable}
\ifLuaTeX
\usepackage[bidi=basic,shorthands=off,provide=*]{babel}
\else
\usepackage[bidi=default,shorthands=off,provide=*]{babel}
\fi
\ifPDFTeX
\else
\babelfont{rm}[]{Times New Roman}
\fi
\ifLuaTeX
  \usepackage{selnolig} % disable illegal ligatures
\fi
\setlength{\emergencystretch}{3em} % prevent overfull lines
\providecommand{\tightlist}{%
  \setlength{\itemsep}{0pt}\setlength{\parskip}{0pt}}
\usepackage{microtype}
\usepackage{xurl}
\usepackage{etoolbox}
\setlength{\emergencystretch}{3em}
\setlength{\tabcolsep}{4pt}
\renewcommand{\arraystretch}{1.15}
\AtBeginEnvironment{longtable}{\small}
\usepackage{bookmark}
\IfFileExists{xurl.sty}{\usepackage{xurl}}{} % add URL line breaks if available
\urlstyle{same}
% fallback for those not using the hyperref driver hyperxmp:
\makeatletter
\@ifundefined{xmpquote}{\newcommand{\xmpquote}[1]{#1}}{}
\makeatother
\hypersetup{
  pdftitle={领域 roadmap：随机快速不动点与单调算子方法},
  pdfauthor={Astra 调研版},
  pdflang={zh-CN},
  colorlinks=true,
  linkcolor={Maroon},
  filecolor={Maroon},
  citecolor={Blue},
  urlcolor={Blue},
  pdfcreator={LaTeX via pandoc}}

\title{领域 roadmap：随机快速不动点与单调算子方法}
\author{Astra 调研版}
\date{2026-10-14}

\begin{document}
\maketitle

\renewcommand*\contentsname{目录}
{
\setcounter{tocdepth}{2}
\tableofcontents
}
\subsection{1. 课题判断}\label{ux8bfeux9898ux5224ux65ad}

\textbf{``随机 Fast KM 与随机 Fast OGDA
的几乎必然快速率''已有直接近邻，三重交集不能整体标为尚无结果。}
VR-FKM、AFP、VFOSA 和 VAPEG 已在各自假设下证明末次平方残差的
\(o(k^{-2})\)
几乎必然速率，并证明迭代点几乎必然趋于解。最接近用户目标的证据为
{[}R16{]} Theorem 2、{[}R17{]} Theorem 2、{[}R18{]} Theorem 14 和
{[}R19{]} Theorem 3.2。

\textbf{值得继续的主线是：冻结 Boţ
系列原始递推，判明其随机过程是否已被上述框架覆盖，再研究有实质差异的噪声门槛。}
仅加入随机 oracle、递增 batch，或把期望分析换成 Robbins--Siegmund
引理，均不足以确认新颖性。原始 Fast KM 的旧噪声缓存、Hilbert
空间弱点极限、原始 Fast OGDA 的查询点与误差传播，是需要逐式比较的对象。

本文保留任务书的四个交集，但将 B 理解为``加速机制''，包括 Nesterov
动量、Halpern 锚点和 optimistic / past-extragradient
机制。它们有联系，但名称和更新式不能互换。有限时域的随机复杂度列在
B∩C，只有明确含同一无限轨道 a.s. 结论的随机加速结果列入 A∩B∩C。

本报告覆盖本地两份讨论材料、Zotero 的 fixed-point PDF、原论文正文、arXiv
版本记录、期刊或会议页面，以及研究者主页。检索截至
\textbf{2026-10-14}。``本次核实的最强界''仅指所列假设和输出指标下的结果，不代表全部
oracle 模型的统一最优界。

\subsubsection{1.1
读表之前必须固定的指标}\label{ux8bfbux8868ux4e4bux524dux5fc5ux987bux56faux5b9aux7684ux6307ux6807}

对不动点问题取 \(V=I-T\)，记 \(r_k=\lVert x_k-Tx_k\rVert\)；对求零问题记
\(r_k=\lVert Vx_k\rVert\)。表内 \(r_k\) 总是\textbf{范数}，\(r_k^2\)
是平方。优化目标差另记
\(h_k=f(x_k)-\min f\)。不同算法查询点不同时，单独注明。

\begin{itemize}
\tightlist
\item
  \textbf{a.s.} 是 almost surely，即概率为一的轨道性质。
\item
  \textbf{末次迭代} 指确定索引 \(k\) 的算法输出；running minimum
  是此前最小残差；随机输出按另一个分布选取索引。三者不互换。
\item
  \textbf{RMS 残差} 为 \((\mathbb E r_k^2)^{1/2}\)。它控制
  \(\mathbb E r_k\)，反方向不能直接推得。
\item
  \textbf{oracle} 是一次随机梯度或随机算子查询。\(k\) 次迭代若每次用
  \(b_k\) 个样本，总成本为 \(N_k=\sum_{j\le k}b_j\)。
\item
  \(\widetilde O\) 隐去对数因子；复杂度目标为范数不超过
  \(\varepsilon\)，平方残差目标相应为 \(\varepsilon^2\)。
\item
  \textbf{弱点极限}
  是所有连续线性泛函意义下的点极限。有限维中弱、强点极限一致，无限维中需区分。
\end{itemize}

仅有 \(\mathbb E r_k^2=O(k^{-2})\) 不能直接推出 \(r_k=o(k^{-1})\)
a.s.。同样，为每个终止时刻 \(K\)
重新选择参数的定理，不能直接解释为一条无限轨道的定理。

\subsubsection{1.2
问题类与噪声}\label{ux95eeux9898ux7c7bux4e0eux566aux58f0}

\textbf{非扩张映射}满足 \(\lVert Tx-Ty\rVert\le\lVert x-y\rVert\)。在
Hilbert 空间中，\(V=I-T\) 为
\(1/2\)-余强制算子。\textbf{\(\beta\)-余强制}指

\[
\langle Vx-Vy,x-y\rangle\ge\beta\lVert Vx-Vy\rVert^2.
\]

余强制蕴含单调与 Lipschitz 连续；反方向对一般算子不成立。因此，Fast KM
的问题类包含于单调 Lipschitz 求零问题类，但 Fast KM 的递推不等于任意
Fast OGDA 特例。{[}R10{]} Remark 2.2 已明确讨论两者更新式的差异。

下文用 N2
表示条件无偏且二阶误差有界：\(\mathbb E[e_k\mid\mathcal F_k]=0\)、\(\mathbb E[\lVert e_k\rVert^2\mid\mathcal F_k]\le\sigma_k^2\)。若只要求
\(\sigma_k\le\sigma\)，噪声不必消失。N6 必须同时写明权重，例如
\(\sum k^2\sigma_k^2<\infty\)；单独说``噪声平方可和''没有足够信息。SGC
是 strong growth
condition，指随机梯度的二阶矩由真实梯度的平方控制，因而噪声随最优性残差消失。

\subsection{2. 四张速率总表}\label{ux56dbux5f20ux901fux7387ux603bux8868}

\subsubsection{2.1
A∩B：光滑凸优化中的随机动量}\label{abux5149ux6ed1ux51f8ux4f18ux5316ux4e2dux7684ux968fux673aux52a8ux91cf}

本表均考察光滑凸目标；Salzo
的复合目标定理在正则项取零时属于该类。\textbf{阻尼系数 \(1-\alpha/k\)
与梯度步长 \(\alpha/k\) 是两个不同参数。}
此处先记录原文能确认的参数族，不把所有随机 Nesterov 变体归于一个速率。

\textbf{条件表。} 与下方速率表按方法名称对应。

{\def\LTcaptype{none} % do not increment counter
\begin{longtable}[]{@{}
  >{\raggedright\arraybackslash}p{(\linewidth - 4\tabcolsep) * \real{0.2400}}
  >{\raggedright\arraybackslash}p{(\linewidth - 4\tabcolsep) * \real{0.2800}}
  >{\raggedright\arraybackslash}p{(\linewidth - 4\tabcolsep) * \real{0.4800}}@{}}
\toprule\noalign{}
\begin{minipage}[b]{\linewidth}\raggedright
方法
\end{minipage} & \begin{minipage}[b]{\linewidth}\raggedright
问题类
\end{minipage} & \begin{minipage}[b]{\linewidth}\raggedright
噪声与参数
\end{minipage} \\
\midrule\noalign{}
\endhead
\bottomrule\noalign{}
\endlastfoot
SNAG，固定动量 & 光滑凸，有限维 & N2/原文 ABC 二阶矩；步长
\(\Theta(k^{-2/3-\delta})\) \\
SNAG，双步长、多序列 & 光滑凸有限和 & SGC；按定理选两组参数 \\
SNAG，\(1-\alpha/k\) 阻尼 & 光滑凸，Hilbert &
常梯度步长；\(\alpha>3\)；\(\sum k^2\sigma_k^2<\infty\) \\
同一 SNAG，点极限 & 光滑凸，可分 Hilbert & 上行并加强为
\(\sum k\sigma_k<\infty\) \\
stochastic Ravine & 光滑凸，可分 Hilbert &
\(\alpha>3\)；常步长；\(\sum k\sigma_k<\infty\) \\
上述 SNAG，递增批量 & 光滑凸 &
独立样本方差有界；\(b_k\asymp k^{3+\delta}\) \\
stochastic FISTA & 光滑凸特例，Hilbert &
有界条件方差；\(t_k\asymp k\)；\(\gamma_k\asymp k^{-3/2}(\log k)^{-r}\)，\(r>1/2\) \\
SHB-IMA，参照线 & 光滑凸随机优化 &
原文平滑/二阶矩条件；\(\eta_k\asymp k^{-1/2-\delta}\) \\
\end{longtable}
}

\textbf{速率表。}

{\def\LTcaptype{none} % do not increment counter
\begin{longtable}[]{@{}
  >{\raggedright\arraybackslash}p{(\linewidth - 6\tabcolsep) * \real{0.2200}}
  >{\raggedright\arraybackslash}p{(\linewidth - 6\tabcolsep) * \real{0.2100}}
  >{\raggedright\arraybackslash}p{(\linewidth - 6\tabcolsep) * \real{0.3500}}
  >{\raggedright\arraybackslash}p{(\linewidth - 6\tabcolsep) * \real{0.2200}}@{}}
\toprule\noalign{}
\begin{minipage}[b]{\linewidth}\raggedright
方法
\end{minipage} & \begin{minipage}[b]{\linewidth}\raggedright
结论类型
\end{minipage} & \begin{minipage}[b]{\linewidth}\raggedright
速率或极限
\end{minipage} & \begin{minipage}[b]{\linewidth}\raggedright
论文与定理
\end{minipage} \\
\midrule\noalign{}
\endhead
\bottomrule\noalign{}
\endlastfoot
SNAG，固定动量 & 末次 a.s.；点 a.s. 极限 &
\(h_k=O(k^{-1/3+\delta})\)，\(0<\delta<1/3\) & {[}R3{]} Thm. 5 \\
SNAG，双步长、多序列 & 末次期望；末次 a.s. &
\(\mathbb E h_k=O(k^{-2})\)；\(h_k=o(k^{-2})\) a.s. & {[}R4{]} Thms.
3--4 \\
SNAG，\(1-\alpha/k\) 阻尼 & 末次 a.s. & \(h_k=o(k^{-2})\) & {[}R5{]}
Thm. 3.1(ii) 专门化 \\
同一 SNAG，点极限 & 点 a.s. 弱极限 & 点趋于最优解；目标差仍
\(o(k^{-2})\) & {[}R5{]} Thm. 3.1(iii) \\
stochastic Ravine & 末次 a.s.；点 a.s. 弱极限 & \(h_k=o(k^{-2})\) &
{[}R5{]} Cor. 4.1 \\
上述 SNAG，递增批量 & 末次 a.s. 目标差 &
\(h_k=o(k^{-2})\)；\(N_k=\Theta(k^{4+\delta})\) & {[}R5{]} Thm. 3.1(ii)
的批量推论 \\
stochastic FISTA & 末次期望；点 a.s. 弱极限 &
\(\mathbb E h_k=O((\log k)^r/\sqrt{k})\) & {[}R6{]} Thm. 3.10、Rem.
3.11 \\
SHB-IMA，参照线 & 末次 a.s.；点 a.s. 极限 & \(h_k=o(k^{-1/2+\delta})\) &
{[}R2{]} Thm. 13、Cor. 17 \\
\end{longtable}
}

\textbf{条件说明。} {[}R5{]} 的一般权重为
\(\sum s_k^2t_k^2\sigma_k^2<\infty\)，大 O 到小 o
还需要加强的外推参数条件；上表取
\(t_k\asymp k\)、常步长。点极限定理采用更强的加权一阶可和条件，不能只凭目标值速率补上点极限。对于批量实现，\(b_k\asymp k^{3+\delta}\)
保证前一平方可和条件；若套用 \(\sum k\sigma_k<\infty\)，足够的幂次是
\(b_k\asymp k^{4+\delta}\)。

{[}R4{]} 使用 SGC，不能宣称持续加性方差下仍有相同结果。{[}R2{]} 是
heavy-ball 参照，不把它记为 Nesterov 定理。Salzo
的结果支持无限轨道的点极限，但表中列出的速率是期望目标差，不能写成 a.s.
\(o(k^{-2})\)。

\textbf{梯度步长恰为 \(\alpha/k\) 的格子：空白，待核实。} 本次核查
{[}R2{]}--{[}R6{]}、本地随机 Nesterov 笔记及 arXiv
线索后，未确认适用于``任意光滑凸目标、持续 N2、原始指定 Nesterov
递推''的统一多项式末次 a.s.
速率。该标记仅表示本报告未确认，不能据此声称不存在研究。

\subsubsection{2.2 A∩C：随机 KM、Halpern
与轨道理论}\label{acux968fux673a-kmhalpern-ux4e0eux8f68ux9053ux7406ux8bba}

\textbf{条件表。} 与下方速率表按方法名称对应。

{\def\LTcaptype{none} % do not increment counter
\begin{longtable}[]{@{}
  >{\raggedright\arraybackslash}p{(\linewidth - 4\tabcolsep) * \real{0.2400}}
  >{\raggedright\arraybackslash}p{(\linewidth - 4\tabcolsep) * \real{0.2800}}
  >{\raggedright\arraybackslash}p{(\linewidth - 4\tabcolsep) * \real{0.4800}}@{}}
\toprule\noalign{}
\begin{minipage}[b]{\linewidth}\raggedright
方法
\end{minipage} & \begin{minipage}[b]{\linewidth}\raggedright
问题类
\end{minipage} & \begin{minipage}[b]{\linewidth}\raggedright
噪声与参数
\end{minipage} \\
\midrule\noalign{}
\endhead
\bottomrule\noalign{}
\endlastfoot
stochastic KM & 非扩张；Hilbert &
\(\sum\lambda_k(\mathbb E_k\lVert e_k\rVert^2)^{1/2}<\infty\)；\(\sum\lambda_k(1-\lambda_k)=\infty\) \\
stochastic KM & 非扩张；有限维一般范数 & 鞅差；原文噪声、步长条件 \\
stochastic KM & 同上 & 有界方差；\(\lambda_k\asymp k^{-2/3}\) \\
stochastic KM & Euclidean 非扩张 & 已知 \(K\)；常步长约 \(K^{-1/2}\) \\
stochastic KM & 可分 Hilbert；随机映射非扩张等原文结构 &
仅一个解处有限方差；RM 步长 \\
同上，running minimum & 同上 &
\(\lambda_k\asymp k^{-a}\)，\(1/2<a<1\) \\
同上，末次期望 & 同上 & \(a=2/3\)；原文矩条件 \\
stochastic Halpern & 非扩张；有限维一般范数 &
累积加权误差受控；递增批量 \\
stochastic Halpern & 同上；点极限另需光滑范数 &
\(\sum\mathbb E\lVert U_k\rVert<\infty\) \\
\end{longtable}
}

\textbf{速率表。}

{\def\LTcaptype{none} % do not increment counter
\begin{longtable}[]{@{}
  >{\raggedright\arraybackslash}p{(\linewidth - 6\tabcolsep) * \real{0.2200}}
  >{\raggedright\arraybackslash}p{(\linewidth - 6\tabcolsep) * \real{0.2100}}
  >{\raggedright\arraybackslash}p{(\linewidth - 6\tabcolsep) * \real{0.3500}}
  >{\raggedright\arraybackslash}p{(\linewidth - 6\tabcolsep) * \real{0.2200}}@{}}
\toprule\noalign{}
\begin{minipage}[b]{\linewidth}\raggedright
方法
\end{minipage} & \begin{minipage}[b]{\linewidth}\raggedright
结论类型
\end{minipage} & \begin{minipage}[b]{\linewidth}\raggedright
速率或极限
\end{minipage} & \begin{minipage}[b]{\linewidth}\raggedright
论文与定理
\end{minipage} \\
\midrule\noalign{}
\endhead
\bottomrule\noalign{}
\endlastfoot
stochastic KM & 点 a.s. 弱极限 & \(x_k\) 趋于固定点；不提供通用数值阶 &
{[}R1{]} Cor. 2.7 \\
stochastic KM & 点 a.s. 极限 & 固定点极限 & {[}R7{]} Thms. 2.5、2.8 \\
stochastic KM & 末次期望范数 & \(\mathbb E r_k=O(\log k/k^{1/6})\) &
{[}R7{]} Thm. 4.4 \\
stochastic KM & 随机输出期望范数 &
\(O(K^{-1/4})\)；\(O(\varepsilon^{-4})\) 样本 & {[}R7{]} Cor. 5.4 \\
stochastic KM & 末次残差 a.s.；点 a.s. 弱极限 & \(r_k\to0\) & {[}R8{]}
Thm. 2.13 \\
同上，running minimum & 最佳历史点 a.s. &
\(\min_{j\le k}r_j=o(k^{-(1-a)/2})\) & {[}R8{]} Cor. 2.12 \\
同上，末次期望 & 末次期望范数 & \(O(\sqrt{\log k}/k^{1/6})\) & {[}R8{]}
Thm. 2.11、Cor. 2.12 \\
stochastic Halpern & 末次期望范数 & \(\widetilde O(k^{-1})\)；批量方案
\(\widetilde O(\varepsilon^{-5})\) & {[}R9{]} Thm. 3.3、Cor. 3.5 \\
stochastic Halpern & 残差 a.s.；点 a.s. 极限 & \(r_k\to0\)；该定理未给
\(o(k^{-1})\) a.s. & {[}R9{]} Thm. 3.10 \\
\end{longtable}
}

{[}R8{]} Corollary 2.12 同时列出随机输出、running minimum
与末次期望界；表中按三种不同输出分列。{[}R1{]} Corollary 2.7
允许有偏误差，但采用加权条件标准差可和；不能把它写成仅需条件方差有界。{[}R8{]}
也不是仅把任意 oracle
的统一方差界替换成``解处方差有限''：随机映射的结构假设仍须满足。

Halpern
属于锚点加速。本表强调其轨道极限，下一表补充它的确定性速率及有限精度随机复杂度。因此，同一种方法可以出现在不同交集，但每行证明的性质不同。

\subsubsection{2.3
B∩C：确定性快速率与随机样本复杂度}\label{bcux786eux5b9aux6027ux5febux901fux7387ux4e0eux968fux673aux6837ux672cux590dux6742ux5ea6}

\textbf{条件表。} 与下方速率表按方法名称对应。

{\def\LTcaptype{none} % do not increment counter
\begin{longtable}[]{@{}
  >{\raggedright\arraybackslash}p{(\linewidth - 4\tabcolsep) * \real{0.2400}}
  >{\raggedright\arraybackslash}p{(\linewidth - 4\tabcolsep) * \real{0.2800}}
  >{\raggedright\arraybackslash}p{(\linewidth - 4\tabcolsep) * \real{0.4800}}@{}}
\toprule\noalign{}
\begin{minipage}[b]{\linewidth}\raggedright
方法
\end{minipage} & \begin{minipage}[b]{\linewidth}\raggedright
问题类
\end{minipage} & \begin{minipage}[b]{\linewidth}\raggedright
噪声与参数
\end{minipage} \\
\midrule\noalign{}
\endhead
\bottomrule\noalign{}
\endlastfoot
KM & 非扩张；有界轨道；一般范数 & 无噪声；松弛远离 0、1 \\
Halpern & 非扩张；Hilbert & 无噪声；经典锚点权重 \\
Fast KM & \(\theta\)-averaged，含 \(\theta=1\) 非扩张；Hilbert &
无噪声；\(\alpha>2\)，\(0<s\le1/\theta\) \\
generalized Fast KM & 非扩张；Hilbert & 无噪声；原文 overrelaxation
参数 \\
EAG-C / EAG-V & 单调 Lipschitz；有限维 & 无噪声；原文步长 \\
G-EAG & 单调 Lipschitz；Hilbert & 无噪声；衰减 Tikhonov 锚点 \\
Fast OGDA & 单调 \(L\)-Lipschitz；Hilbert &
无噪声；\(\alpha>2\)，\(0<s<1/(2L)\) \\
VR-Halpern & 余强制及单调 Lipschitz & 无偏方差界；同样本多点差分、VR
结构 \\
stochastic iKM & 非扩张；实 Hilbert 空间 & N2，可加受控有偏项；预设
\(K\) \\
stochastic Dual-OHM & 期望余强制 / 平方非扩张 & 常批量 \(B\)；预设
\(N\) \\
dual-anchor + 内层近似 & 单调 Lipschitz 求零 &
样本算子单调、Lipschitz；N2 \\
anytime stochastic anchoring & 单调 Lipschitz；约束问题亦可 &
一次新样本；原文二阶矩结构 \\
同上，VR 版本 & 同上 & 每步两个样本；同样本多点访问 \\
VR-GHAL & 非扩张；适当 Banach 几何 & 样本映射非扩张；N2 \\
\end{longtable}
}

\textbf{速率表。}

{\def\LTcaptype{none} % do not increment counter
\begin{longtable}[]{@{}
  >{\raggedright\arraybackslash}p{(\linewidth - 6\tabcolsep) * \real{0.2200}}
  >{\raggedright\arraybackslash}p{(\linewidth - 6\tabcolsep) * \real{0.2100}}
  >{\raggedright\arraybackslash}p{(\linewidth - 6\tabcolsep) * \real{0.3500}}
  >{\raggedright\arraybackslash}p{(\linewidth - 6\tabcolsep) * \real{0.2200}}@{}}
\toprule\noalign{}
\begin{minipage}[b]{\linewidth}\raggedright
方法
\end{minipage} & \begin{minipage}[b]{\linewidth}\raggedright
结论类型
\end{minipage} & \begin{minipage}[b]{\linewidth}\raggedright
速率或极限
\end{minipage} & \begin{minipage}[b]{\linewidth}\raggedright
论文与定理
\end{minipage} \\
\midrule\noalign{}
\endhead
\bottomrule\noalign{}
\endlastfoot
KM & 确定性末次范数 & \(O(k^{-1/2})\) & {[}R11{]} Thm. 1.1 无误差特例 \\
Halpern & 确定性末次范数 & \(r_k\le2\lVert x_0-x_*\rVert/(k+1)\) &
{[}R12{]} Thm. 2.1 \\
Fast KM & 确定性末次；弱点极限 & \(r_k=o(k^{-1})\) & {[}R10{]} Thms.
3.4--3.5 \\
generalized Fast KM & 确定性末次；弱点极限 & \(\alpha=2\) 为
\(O(k^{-1})\)；\(\alpha>2\) 为 \(o(k^{-1})\) & {[}R15{]} Thm. 2 \\
EAG-C / EAG-V & 确定性末次 & \(r_k^2=O(k^{-2})\) & {[}R13{]} Thms.
1--2；§2.3 的算子推广 \\
G-EAG & 确定性末次；强点极限 & \(r_k=O(k^{-1})\)；最小范数解 & {[}R14{]}
Thm. 2.4、Cor. 2.6 \\
Fast OGDA & 确定性末次；弱点极限 & \(r_k=o(k^{-1})\) & {[}R20{]} Thms.
16--17 \\
VR-Halpern & 有限精度末次期望 &
\(\mathbb E r_k\le C/k+C'\varepsilon\)；\(O(\varepsilon^{-3})\) oracle &
{[}R21{]} Thm. 3.1、Cor. 3.4、Thm. 4.1 \\
stochastic iKM & 末次 RMS；样本复杂度 & 无偏时
\(O(K^{-1}+\sigma\log K/\sqrt K)\)；\(\widetilde O(\varepsilon^{-2})\) &
{[}R22{]} Thm. 1 \\
stochastic Dual-OHM & 末次平方残差期望 &
\(O(N^{-2}+\sigma^2/B)\)；\(O(\varepsilon^{-3})\) oracle & {[}R23{]}
Thm. 4.1 \\
dual-anchor + 内层近似 & 末次期望；样本复杂度 & 近最优
\(\widetilde O(\varepsilon^{-2})\) & {[}R24{]} Cor. 3.6 \\
anytime stochastic anchoring & 末次期望范数 & \(O(k^{-1/4})\) &
{[}R25{]} Thm. 2.1、Cor. 2.3 \\
同上，VR 版本 & 末次期望范数 & \(O(k^{-1/2})\) & {[}R25{]} Thm. 3.2 \\
VR-GHAL & anytime 高概率残差 & \(\widetilde O(\varepsilon^{-2})\) oracle
& {[}R26{]} Thm. 1、Cor. 3 \\
\end{longtable}
}

\textbf{表中若干边界。} {[}R22{]} 每步一次新 oracle
调用，但参数依赖总时域；Appendix A Proposition 1 证明的确定性 Fast KM
等价参数与主定理的参数选择不同。{[}R23{]}
的批量在一次运行内不增加，但随目标精度取
\(B=\Theta(\varepsilon^{-2})\)；不能说固定一个 \(B\)
后无限运行得到零残差。{[}R24{]} 包含内层 resolvent 近似，不能称为原始
Fast OGDA 的一步随机化。{[}R25{]}
在约束情形采用梯度映射或原文残差；只有集合值部分为零时才能直接记成
\(\lVert Vx_k\rVert\)。

{[}R26{]} 的 \(\widetilde O(\varepsilon^{-2})\)
使用样本映射非扩张结构。只有 N2 时，其 Corollary 1
给出一般范数模型下更昂贵的复杂度；不能去掉结构假设后保留相同指数。高概率且对时域一致的定理，也不自动给出原始
Fast KM 的 a.s. 小 o 速率和点极限。

\subsubsection{2.4
A∩B∩C：已有随机快速轨道结果与用户目标格子}\label{abcux5df2ux6709ux968fux673aux5febux901fux8f68ux9053ux7ed3ux679cux4e0eux7528ux6237ux76eeux6807ux683cux5b50}

\textbf{条件表。} 与下方速率表按方法名称对应。

{\def\LTcaptype{none} % do not increment counter
\begin{longtable}[]{@{}
  >{\raggedright\arraybackslash}p{(\linewidth - 4\tabcolsep) * \real{0.2400}}
  >{\raggedright\arraybackslash}p{(\linewidth - 4\tabcolsep) * \real{0.2800}}
  >{\raggedright\arraybackslash}p{(\linewidth - 4\tabcolsep) * \real{0.4800}}@{}}
\toprule\noalign{}
\begin{minipage}[b]{\linewidth}\raggedright
方法
\end{minipage} & \begin{minipage}[b]{\linewidth}\raggedright
问题类
\end{minipage} & \begin{minipage}[b]{\linewidth}\raggedright
噪声与参数
\end{minipage} \\
\midrule\noalign{}
\endhead
\bottomrule\noalign{}
\endlastfoot
VR-FKM & 有限和余强制型；有限维 & 分量平均余强制型条件；SVRG / SAGA \\
AFP & 余强制 / 非扩张；有限维 &
递推误差条件；允许有偏与延迟；\(\sum\delta_k<\infty\) \\
AFP，一般随机 oracle & 同上 &
N2；\(b_k\gtrsim k^{3+\omega}\)，\(\omega>0\) \\
VFOSA & 单调复合包含的原文特例；有限维 & 余强制型结构、VR
误差递推及可和尾项 \\
VAPEG，旧称 VFOG & 单调 Lipschitz 求零特例；有限维 & 一般误差条件
(15)；可和尾项；递增 batch 可实现 \\
stochastic Fast OGDA 动力系统 & 单调算子；Hilbert &
连续时间扩散；\(\int s^2\lVert\sigma_Z(s)\rVert_{HS}^2ds<\infty\) \\
\textbf{用户目标：原始 Fast KM} & 非扩张 / 余强制；Hilbert & 缓存旧
oracle；先 N6，再判定 N2 下可行目标 \\
\textbf{用户目标：原始 Fast OGDA} & 单调 Lipschitz；Hilbert &
原始查询点、缓存方式及噪声模型必须固定 \\
\end{longtable}
}

\textbf{速率表。}

{\def\LTcaptype{none} % do not increment counter
\begin{longtable}[]{@{}
  >{\raggedright\arraybackslash}p{(\linewidth - 6\tabcolsep) * \real{0.2200}}
  >{\raggedright\arraybackslash}p{(\linewidth - 6\tabcolsep) * \real{0.2100}}
  >{\raggedright\arraybackslash}p{(\linewidth - 6\tabcolsep) * \real{0.3500}}
  >{\raggedright\arraybackslash}p{(\linewidth - 6\tabcolsep) * \real{0.2200}}@{}}
\toprule\noalign{}
\begin{minipage}[b]{\linewidth}\raggedright
方法
\end{minipage} & \begin{minipage}[b]{\linewidth}\raggedright
结论类型
\end{minipage} & \begin{minipage}[b]{\linewidth}\raggedright
速率或极限
\end{minipage} & \begin{minipage}[b]{\linewidth}\raggedright
论文与定理
\end{minipage} \\
\midrule\noalign{}
\endhead
\bottomrule\noalign{}
\endlastfoot
VR-FKM & 末次期望；末次 a.s.；点 a.s. 极限 &
\(\mathbb E r_k^2=o(k^{-2})\)；\(r_k=o(k^{-1})\) a.s. & {[}R16{]} Thms.
1--2；Cors. 1--2 \\
AFP & 查询点残差 a.s.；点 a.s. 极限 & \(\lVert Gy_k\rVert=o(k^{-1})\)
a.s. & {[}R17{]} Thm. 2 \\
AFP，一般随机 oracle & a.s. 快速率；期望复杂度 & 范数 \(o(k^{-1})\)
a.s.；\(O(\varepsilon^{-4-\omega})\) oracle & {[}R17{]} Thm. 2、Cor.
2 \\
VFOSA & FBS 残差 a.s.；点 a.s. 极限 &
\(\lVert\mathcal G_\lambda x_k\rVert=o(k^{-1})\) a.s. & {[}R18{]} Thms.
12--14 \\
VAPEG，旧称 VFOG & 末次残差 a.s.；点 a.s. 极限 &
\(\lVert Gx_k\rVert=o(k^{-1})\) a.s. & {[}R19{]} Thm. 3.2；Cor. 3.4 \\
stochastic Fast OGDA 动力系统 & 轨道 a.s.；弱点极限 &
\(\lVert VZ(s)\rVert=o(s^{-1})\) a.s. & {[}R27{]} Thm. 3.5 \\
\textbf{用户目标：原始 Fast KM} & 同一无限轨道；末次 a.s. 与弱点极限 &
目标 \(o(k^{-1})\)；\textbf{宽泛结论已有，原始随机递推覆盖待验证} & 对比
{[}R10{]}、{[}R16{]}、{[}R17{]}、{[}R22{]} \\
\textbf{用户目标：原始 Fast OGDA} & 同一无限轨道；末次 a.s. 与弱点极限 &
目标 \(o(k^{-1})\)；\textbf{宽泛结论已有，原始随机递推覆盖待验证} & 对比
{[}R19{]}、{[}R20{]}、{[}R27{]} \\
\end{longtable}
}

VFOSA 的残差是 forward-backward splitting 残差。它包含
resolvent，不能对一般复合包含直接替换成 \(\lVert Vx_k\rVert\)。VAPEG
取集合值算子为零、单调参数为零，即得到本任务允许的单调 Lipschitz
求零特例；本文不展开其非单调部分。

AFP 的``查询点'' \(y_k\)
具有明示快速率，三个序列的点极限另由定理给出。不要只因为 \(x_k\) 与
\(y_k\) 有同一极限，就自动把 \(o(k^{-1})\) 残差率转移到 \(x_k\)。

\textbf{本交集未填写无文献依据的``空白''。}
真正未解决的是覆盖判定及假设边界。本次已查本地笔记、Zotero、arXiv、期刊正文、作者主页和这些文献的相关工作部分；该范围仍不构成不存在其他定理的证明。

\subsubsection{2.5
下界与比较规则}\label{ux4e0bux754cux4e0eux6bd4ux8f83ux89c4ux5219}

{\def\LTcaptype{none} % do not increment counter
\begin{longtable}[]{@{}
  >{\raggedright\arraybackslash}p{(\linewidth - 6\tabcolsep) * \real{0.2200}}
  >{\raggedright\arraybackslash}p{(\linewidth - 6\tabcolsep) * \real{0.2100}}
  >{\raggedright\arraybackslash}p{(\linewidth - 6\tabcolsep) * \real{0.3500}}
  >{\raggedright\arraybackslash}p{(\linewidth - 6\tabcolsep) * \real{0.2200}}@{}}
\toprule\noalign{}
\begin{minipage}[b]{\linewidth}\raggedright
下界
\end{minipage} & \begin{minipage}[b]{\linewidth}\raggedright
适用模型
\end{minipage} & \begin{minipage}[b]{\linewidth}\raggedright
含义与限制
\end{minipage} & \begin{minipage}[b]{\linewidth}\raggedright
证据
\end{minipage} \\
\midrule\noalign{}
\endhead
\bottomrule\noalign{}
\endlastfoot
范数残差 \(\Omega(1/K)\) & 确定性非扩张不动点；原文 oracle 类与维度条件
& Halpern 的最坏情形阶已最优 & {[}R28{]} Thm. 4.6 \\
平方残差 \(\Omega(1/K^2)\) & 光滑凸凹、相应一阶 oracle 类 & 与 EAG 的
\(O(K^{-2})\) 平方率配对 & {[}R13{]} Thm. 3、Cor. 3 \\
随机样本 \(\Omega(\sigma^2\varepsilon^{-2})\) & 含带噪均值估计的 Hilbert
求零子类 & 仅说明噪声项的精度指数；不统一覆盖一般范数 & {[}R22{]}
下界讨论；简约子类见下文 \\
随机样本 \(\Omega(\varepsilon^{-3})\) & {[}R9{]}
规定的平均化算法、batching 和一般范数模型 & 不能用于否定 Hilbert 中更强
oracle 的二次方复杂度 & {[}R9{]} Thm. 3.7 \\
\end{longtable}
}

第一、二行的最坏情形允许困难实例随 \(K\) 改变。Fast KM 的 \(o(1/k)\)
是每个固定实例沿无限轨道的渐近结论，两者没有逻辑冲突。

第三行可由 \(V_\mu(x)=x-\mu\)、oracle \(V_\mu(x)+\xi\)
的高斯均值估计子问题解释：一次查询等价于获得一个关于 \(\mu\)
的带噪样本；要令均方残差不超过 \(\varepsilon^2\)，至少需要常数倍
\(\sigma^2/\varepsilon^2\)
个样本。这里是对子类的标准下界归约，未将它冒充 {[}R22{]} 的有编号定理。

2026 年还存在 {[}R29{]} 的 Hilbert 归约与 {[}R30{]} 的一般范数 oracle
复杂度研究。它们提示\textbf{空间几何与 oracle
访问方式会改变下界}。{[}R30{]}
本次已确认元数据，未完成正文定理核验；其数值界不纳入总表。

\subsection{3.
直接相关的发展时间线}\label{ux76f4ux63a5ux76f8ux5173ux7684ux53d1ux5c55ux65f6ux95f4ux7ebf}

以下 18
个节点按首次相关成果出现的年份组织；预印本年份与正式出版年份不同的，均予以注明。

{\def\LTcaptype{none} % do not increment counter
\begin{longtable}[]{@{}
  >{\raggedright\arraybackslash}p{(\linewidth - 4\tabcolsep) * \real{0.1200}}
  >{\raggedright\arraybackslash}p{(\linewidth - 4\tabcolsep) * \real{0.3600}}
  >{\raggedright\arraybackslash}p{(\linewidth - 4\tabcolsep) * \real{0.5200}}@{}}
\toprule\noalign{}
\begin{minipage}[b]{\linewidth}\raggedright
年份
\end{minipage} & \begin{minipage}[b]{\linewidth}\raggedright
作者与论文
\end{minipage} & \begin{minipage}[b]{\linewidth}\raggedright
推进的结果及总表位置
\end{minipage} \\
\midrule\noalign{}
\endhead
\bottomrule\noalign{}
\endlastfoot
2015 & Combettes--Pesquet，{[}R1{]} & 随机拟 Fejér 框架；随机 KM 的 a.s.
弱点极限，A∩C 的基础 \\
2015 & Boţ--Csetnek--Hendrich，{[}R35{]} & 惯性不动点与
Douglas--Rachford 结构，为 B∩C 提供递推与点极限工具 \\
2017 & Bravo--Cominetti--Pavez-Signé，{[}R11{]} & 不精确 KM
的显式残差误差界，连接误差累积与 A∩C \\
2020--2021 & Lieder，{[}R12{]} & Halpern 的精确 \(O(k^{-1})\)
范数界；2020 online、2021 卷期 \\
2021 & Sebbouh--Gower--Defazio，{[}R2{]} & SHB 末次 a.s.
目标差与点极限，A∩B 的参照线 \\
2021 & Yoon--Ryu，{[}R13{]} & EAG 把单调 Lipschitz 末次平方残差推进到
\(O(k^{-2})\)，B∩C \\
2022 & Liu--Yuan，{[}R3{]} & SGD、SHB、SNAG 的统一 a.s.
分析；一般凸末次目标差结果，A∩B \\
2022 & Park--Ryu，{[}R28{]} & 非扩张不动点的加速复杂度和匹配下界，B∩C \\
2022 & Cai--Song--Guzmán--Diakonikolas，{[}R21{]} & 随机 Halpern
结合递归 VR，有限精度 oracle 成本 \(O(\varepsilon^{-3})\)，B∩C \\
2022--2023 & Boţ--Nguyen，{[}R10{]} & Fast KM；2022 预印本、2023
SINA，\(o(k^{-1})\) 残差及弱点极限 \\
2022--2025 & Boţ--Csetnek--Nguyen，{[}R20{]} & Fast OGDA；2023
online、2025 FoCM 卷期，单调 Lipschitz 的 \(o(k^{-1})\) \\
2024 & Bravo--Cominetti，{[}R7{]} & 随机 KM 的 a.s.
点极限与可比较的期望误差界，A∩C \\
2024--2026 & Tran-Dinh，{[}R16{]} & VR-FKM；本报告使用 2026 v4，其中含
a.s. 小 o 与点极限，三重交集 \\
2025 & Hermant 等，{[}R4{]} & SGC 下随机 Nesterov 的加速及 a.s. 小 o
目标差，A∩B \\
2025--2026 & Tran-Dinh、Nguyen-Trung，{[}R17{]}--{[}R19{]} &
AFP、VFOSA、VAPEG 把随机快速轨道理论推进到更广 oracle 和算子结构 \\
2025--2026 & Salzo，{[}R6{]}；Attouch--Fadili--Kungurtsev，{[}R5{]} &
FISTA 的随机点极限；Ravine 的误差权重、快速目标差与点极限，A∩B \\
2026 &
Boţ--Chenchene，{[}R14{]}；Boţ--Chenchene--Fadili，{[}R15{]}；Boţ--Schindler，{[}R27{]}
& 强点极限、广义 Fast KM、随机连续时间 Fast OGDA，明确离散与连续边界 \\
2026 & Cortild--Cartis--Peypouquet，{[}R8{]}；Yang
等，{[}R22{]}；Yoon--Loizou，{[}R23{]}--{[}R24{]}；Kim--Alacaoglu，{[}R25{]}；Diakonikolas，{[}R26{]}
& 解处方差、近最优样本复杂度、anytime 和高概率结果共同细化前沿；彼此
oracle 模型不相同 \\
\end{longtable}
}

\subsection{4.
研究人员与主页}\label{ux7814ux7a76ux4ebaux5458ux4e0eux4e3bux9875}

\subsubsection{4.1
当前代表研究者}\label{ux5f53ux524dux4ee3ux8868ux7814ux7a76ux8005}

选取 18 位与上述窄领域直接相关、能列出两篇代表作的研究者。文献编号对应第
8 节完整书目。所有下列链接均在 \textbf{2026-10-14
实际访问}。官方个人介绍页也可作为身份入口，但研究组名单不能冒充独立个人主页。

{\def\LTcaptype{none} % do not increment counter
\begin{longtable}[]{@{}
  >{\raggedright\arraybackslash}p{(\linewidth - 8\tabcolsep) * \real{0.1900}}
  >{\raggedright\arraybackslash}p{(\linewidth - 8\tabcolsep) * \real{0.2500}}
  >{\raggedright\arraybackslash}p{(\linewidth - 8\tabcolsep) * \real{0.1800}}
  >{\raggedright\arraybackslash}p{(\linewidth - 8\tabcolsep) * \real{0.2400}}
  >{\raggedright\arraybackslash}p{(\linewidth - 8\tabcolsep) * \real{0.1400}}@{}}
\toprule\noalign{}
\begin{minipage}[b]{\linewidth}\raggedright
姓名
\end{minipage} & \begin{minipage}[b]{\linewidth}\raggedright
机构及核验状态
\end{minipage} & \begin{minipage}[b]{\linewidth}\raggedright
主页
\end{minipage} & \begin{minipage}[b]{\linewidth}\raggedright
代表作两篇
\end{minipage} & \begin{minipage}[b]{\linewidth}\raggedright
研究线
\end{minipage} \\
\midrule\noalign{}
\endhead
\bottomrule\noalign{}
\endlastfoot
Radu Ioan Boţ & University of Vienna &
\href{https://www.mat.univie.ac.at/~rabot/}{个人主页}，已访问 & Fast KM
{[}R10{]}；Fast OGDA {[}R20{]} & B∩C；A∩B∩C \\
Ernö Robert Csetnek & University of Vienna &
\href{https://www.mat.univie.ac.at/~csercser/index.htm}{个人主页}，已访问
& 惯性 DR {[}R35{]}；Fast OGDA {[}R20{]} & B∩C \\
Dang-Khoa Nguyen & 原论文署名 Vienna；当前机构待核实 &
\href{https://www.mat.univie.ac.at/~rabot/group.html}{官方组员记录}，已访问；个人主页待核实
& Fast KM {[}R10{]}；Fast OGDA {[}R20{]} & B∩C \\
Enis Chenchene & 2026 论文署名 Vienna；当前机构待核实 &
\href{https://www.mat.univie.ac.at/~rabot/group.html}{官方组员记录}，已访问；个人主页待核实
& G-EAG {[}R14{]}；广义 Fast KM {[}R15{]} & B∩C \\
Jalal Fadili & ENSICAEN / GREYC &
\href{https://fadili.users.greyc.fr/}{个人主页}，已访问 & stochastic
Ravine {[}R5{]}；广义 Fast KM {[}R15{]} & A∩B；B∩C \\
Quoc Tran-Dinh & UNC Chapel Hill &
\href{https://quoctd.web.unc.edu/}{个人主页}，已访问 & VR-FKM
{[}R16{]}；VFOSA {[}R18{]} & A∩B∩C \\
Nghia Nguyen-Trung & UNC Chapel Hill，PhD candidate &
\href{https://nghiant01.github.io/}{个人主页}，已访问 & AFP
{[}R17{]}；VAPEG {[}R19{]} & A∩B∩C \\
Ernest K. Ryu & UCLA；主页列示休假期间在 OpenAI &
\href{https://ernestryu.com/}{个人主页}，已访问 & EAG
{[}R13{]}；精确最优不动点复杂度 {[}R28{]} & B∩C \\
TaeHo Yoon & Johns Hopkins University，Assistant Research Scientist &
\href{https://tetrzim.github.io/CV/}{个人 CV}，已访问 & EAG
{[}R13{]}；直接随机 dual-anchor {[}R23{]} & B∩C \\
Nicolas Loizou & Johns Hopkins University &
\href{https://nicolasloizou.github.io/index.html}{个人主页}，已访问 &
直接随机 dual-anchor {[}R23{]}；近最优 dual-anchor {[}R24{]} & B∩C \\
Mario Bravo & Universidad de Santiago de Chile &
\href{https://sites.google.com/site/mbravogon/}{个人主页}，已访问 &
stochastic KM {[}R7{]}；stochastic Halpern {[}R9{]} & A∩C \\
Roberto Cominetti & Universidad Adolfo Ibáñez &
\href{https://www.uai.cl/profesores/ingenieria-y-ciencias/roberto-cominetti}{官方个人页}，已访问
& 不精确 KM {[}R11{]}；stochastic KM {[}R7{]} & A∩C \\
Juan Pablo Contreras & Universidad Diego Portales &
\href{https://citylog.udp.cl/persona/juan-pablo-contreras/}{官方个人页}，已访问
& 最优一般范数不动点误差界 {[}R32{]}；stochastic Halpern {[}R9{]} &
A∩C；B∩C \\
Patrick L. Combettes & North Carolina State University &
\href{https://pcombet.math.ncsu.edu/}{个人主页}，已访问 & 随机拟 Fejér I
{[}R1{]}；II {[}R31{]} & A∩C \\
Jean-Christophe Pesquet & CentraleSupélec / Université Paris-Saclay &
\href{https://jc.pesquet.eu/biography/}{主页及履历}，已访问 & 随机拟
Fejér I {[}R1{]}；II {[}R31{]} & A∩C \\
Jelena Diakonikolas & University of Wisconsin--Madison &
\href{https://www.jelena-diakonikolas.com/}{个人主页}，已访问 &
VR-Halpern {[}R21{]}；高概率不动点 {[}R26{]} & A∩C；B∩C \\
Saverio Salzo & Sapienza Università di Roma &
\href{https://www.open.diag.uniroma1.it/users/saverio_salzo}{官方个人页}，已访问
& stochastic FISTA {[}R6{]}；不精确加速 proximal point {[}R33{]} &
A∩B \\
Ahmet Alacaoglu & University of British Columbia &
\href{https://ahmetalacaoglu.github.io/}{个人主页}，已访问 & stochastic
VR for VI {[}R34{]}；anytime 末次界 {[}R25{]} & B∩C \\
\end{longtable}
}

\textbf{人员核验中的修正。} 任务线索里的 ``Baptiste Hermant'' 应为
\textbf{Julien Hermant}。\href{https://j-hermant.github.io/}{个人主页}
于 2026-10-14 已访问，列示 University of Bordeaux 博士生。他与 Marien
Renaud、Jean-François Aujol、Charles Dossal、Aude Rondepierre 的
{[}R4{]} 是 A∩B
的直接入口。本表未为满足``两篇''要求把未核实作品加入其名下。Dang-Khoa
Nguyen 与 Enis Chenchene 在 Boţ 组页列于 former
members；因此不把旧署名当成当前任职。Ernest K. Ryu 的主页已是 UCLA /
OpenAI 信息，不能沿用旧 SNU 任职。

\subsubsection{4.2
经典基础人物}\label{ux7ecfux5178ux57faux7840ux4ebaux7269}

{\def\LTcaptype{none} % do not increment counter
\begin{longtable}[]{@{}
  >{\raggedright\arraybackslash}p{(\linewidth - 2\tabcolsep) * \real{0.5000}}
  >{\raggedright\arraybackslash}p{(\linewidth - 2\tabcolsep) * \real{0.5000}}@{}}
\toprule\noalign{}
\begin{minipage}[b]{\linewidth}\raggedright
人物
\end{minipage} & \begin{minipage}[b]{\linewidth}\raggedright
与本领域直接相关的贡献
\end{minipage} \\
\midrule\noalign{}
\endhead
\bottomrule\noalign{}
\endlastfoot
M. A. Krasnosel'skiĭ、W. R. Mann & 松弛不动点迭代，是 KM
系列的基础；两人合列一行 \\
Benjamin Halpern & 锚点不动点迭代，支撑固定点选择和快速残差研究 \\
Yurii Nesterov & 加速动量及估计序列，支撑 A∩B 和 Fast KM 的方法来源 \\
Boris T. Polyak & heavy-ball 与随机逼近背景；不能与 Nesterov
更新式混同 \\
Herbert Robbins、David Siegmund &
非负近超鞅工具，支撑轨道极限和误差可和分析；两人合列一行 \\
Hedy Attouch & 加速动力系统、消失阻尼、渐近速率与点极限；本报告直接对应
{[}R5{]} \\
\end{longtable}
}

此表共 8 人；只作方法来源说明，不用于填充当前研究者名额。

\subsection{5.
六个可检验的当前课题}\label{ux516dux4e2aux53efux68c0ux9a8cux7684ux5f53ux524dux8bfeux9898}

以下均是\textbf{候选研究问题}。除已引用定理与第 7
节的诊断推导外，本次没有证明新的算法速率。先做第一个课题的覆盖判定，再决定是否进入证明阶段。

\subsubsection{5.1 原始 Fast KM
的随机等价性与误差条件}\label{ux539fux59cb-fast-km-ux7684ux968fux673aux7b49ux4ef7ux6027ux4e0eux8befux5deeux6761ux4ef6}

\textbf{问题。} 在可分 Hilbert 空间中固定 {[}R10{]} 的 Fast KM，取
\(g_k=Vx_k+e_k\)，每步只产生一个新估计并缓存 \(g_{k-1}\)。N2 下要求
\(\sum k^2\sigma_k^2<\infty\)。目标是末次 \(r_k=o(k^{-1})\) a.s.，同时有
a.s. 弱点极限。

\textbf{已知最好近邻。} {[}R16{]} Theorem 2 在有限和结构下给出同阶 a.s.
结果；{[}R17{]} Theorem 2
在有限维一般误差框架下给出同阶查询点残差及点极限。其典型无延迟误差特例已容纳
\(\sum k^2\sigma_k^2<\infty\)。{[}R22{]} Appendix A Proposition 1
给出确定性 Fast KM 与特定 iKM 参数的联系。

\textbf{剩余判定。} 必须核对 AFP
消元后是否得到同一随机递推，包含初值、查询点、缓存误差和全部系数。若完全相同且所需假设已包含，有限维目标属于已有定理的推论。Hilbert
版本还需判断是常规弱极限论证还是实际技术障碍；不能仅凭维度变化确认贡献。

2024--2026 代表文献为
{[}R16{]}、{[}R17{]}、{[}R22{]}。与用户课题关系为\textbf{直接竞争}。适合把它作为第一阶段的覆盖审计与证明复现，不宜预先作为新颖性承诺。

\subsubsection{5.2 在明确尾部假设下减少 Fast KM
的批量增长}\label{ux5728ux660eux786eux5c3eux90e8ux5047ux8bbeux4e0bux51cfux5c11-fast-km-ux7684ux6279ux91cfux589eux957f}

\textbf{问题。} 保持 5.1
的原始递推。将``仅有条件二阶矩''换成明确的轻尾假设，例如有限维条件次高斯误差，或有界样本平均并具有相应浓度界。尝试在
\(\sigma_k^2=O(k^{-q})\)、\(2<q\le3\) 下证明 \(r_k=o(k^{-1})\) a.s.
与点极限。

\textbf{已知最好近邻。} {[}R17{]} Theorem 2、{[}R19{]} Theorem 3.2
的可和二阶误差框架通常可用 \(q>3\) 实现。{[}R26{]}
的高概率工具表明尾部控制与 oracle 结构可以改变复杂度，但其算法并非原始
Fast KM。

\textbf{剩余问题。} 新噪声本身满足 \(k e_k\to0\)
只是必要诊断，不能替代对历史累积的控制。第 7
节说明：仅有二阶矩时，\(2<q\le3\) 的统一 a.s. 小 o
目标会失败。轻尾假设必须写入定理，不能在证明中隐含使用。

2024--2026 代表文献为
{[}R16{]}、{[}R17{]}、{[}R26{]}。关系为\textbf{条件化的直接候选课题}。若对应
\(b_k\asymp k^q\)，总样本为
\(\Theta(K^{q+1})\)；只有证明成立后才能把减少的批量写成贡献。

\subsubsection{5.3 原始 Fast OGDA
的噪声相关性}\label{ux539fux59cb-fast-ogda-ux7684ux566aux58f0ux76f8ux5173ux6027}

\textbf{问题。} 对单调 \(L\)-Lipschitz 算子，按 {[}R20{]} Algorithm 2
固定预测点与更新点，并明确同一样本重用、独立重采样或缓存上一估计的选择。先在
N6 下检验 \(\lVert Vx_k\rVert=o(k^{-1})\) a.s. 与 a.s. 弱点极限。

\textbf{已知最好近邻。} {[}R19{]} Theorem 3.2 已在有限维给出单调
Lipschitz 特例的同阶结论。{[}R27{]} Theorem 3.5
已给出对应连续时间随机动力系统的 a.s. 快速率与弱点极限。

\textbf{剩余问题。} 原始 Fast OGDA 与 VAPEG
是否只是变量重写，必须逐式验证。预测点依赖当前样本时，不能把该点的
oracle
误差直接视为条件无偏。若递推不同，仍须说明这种差异带来了什么可证明好处；不同名称不是理论增量。

2024--2026 代表文献为
{[}R19{]}、{[}R20{]}、{[}R27{]}。关系为\textbf{直接竞争}。建议在 Fast KM
阶段完成后推进，避免两个不确定证明同时展开。

\subsubsection{5.4 持续 N2 下的单一 anytime
惯性轨道}\label{ux6301ux7eed-n2-ux4e0bux7684ux5355ux4e00-anytime-ux60efux6027ux8f68ux9053}

\textbf{问题。}
非扩张不动点，条件无偏且方差有界，不预知终止时刻，每步一个新
oracle。目标先设为 \(r_k\to0\) a.s. 和点极限，再寻找与 oracle
成本兼容的末次幂次率；不预设 \(o(k^{-1})\)。

\textbf{已知最好近邻。} {[}R22{]} Theorem 1 有预设时域的
\(\widetilde O(K^{-1/2})\) RMS 残差。{[}R8{]} Theorem 2.13 给随机 KM
的无限轨道弱点极限。{[}R25{]} Theorem 2.1 有 anytime 末次期望
\(O(k^{-1/4})\)，多点 VR 下 Theorem 3.2 为 \(O(k^{-1/2})\)。

\textbf{剩余问题。} 不能只把固定时域参数中的 \(K\) 替换成 \(k\)。新
schedule 必须重新证明漂移与轨道性质。候选 a.s.
速率应从真实误差权重导出，不能把期望界直接附上 a.s. 标签。

2024--2026 代表文献为
{[}R8{]}、{[}R22{]}、{[}R25{]}。关系为\textbf{可借用技巧的后续方向}，不是原始常步长
Fast KM 的直接推论。

\subsubsection{5.5 从 SGC
到算子相对噪声条件}\label{ux4ece-sgc-ux5230ux7b97ux5b50ux76f8ux5bf9ux566aux58f0ux6761ux4ef6}

\textbf{问题。} 余强制求零，条件无偏误差满足
\(\mathbb E_k\lVert e_k\rVert^2\le c\lVert Vx_k\rVert^2\)。在无需预设递增
batch 的条件下，检验原始 Fast KM 或最小参数调整能否同时达到
\(o(k^{-1})\) a.s. 与点极限。

\textbf{已知最好近邻。} {[}R4{]} Theorem 4 在梯度问题的 SGC 下给出
\(h_k=o(k^{-2})\) a.s.。{[}R16{]} Theorem 2 已用有限和 VR
实现算子快速率；{[}R8{]} 说明解处方差模型值得单独分析。

\textbf{剩余问题。}
梯度有目标函数，普通余强制算子没有该势函数。相对噪声项能否由同一
Lyapunov
的负漂移吸收，是具体证明问题。相对噪声的常数阈值、是否隐含插值结构、以及是否已落入
{[}R16{]} 的估计器类，都必须验证。

2024--2026 代表文献为
{[}R4{]}、{[}R8{]}、{[}R16{]}。关系为\textbf{从用户随机 Nesterov
背景迁移的候选机制}，目前不声明新颖性成立。

\subsubsection{5.6
快速残差与强点极限的兼容性}\label{ux5febux901fux6b8bux5deeux4e0eux5f3aux70b9ux6781ux9650ux7684ux517cux5bb9ux6027}

\textbf{问题。} 非扩张或单调 Lipschitz 问题，使用预定锚点与 N6
噪声，希望 \(r_k=O(k^{-1})\)
a.s.，并强趋于指定投影解。必须说明锚点与噪声权重的精确关系。

\textbf{已知最好近邻。} 确定性 {[}R14{]} Theorem 2.4、Corollary 2.6
实现强点极限和 \(O(k^{-1})\) 残差。{[}R9{]} Theorem 3.10
在有限维光滑范数下给随机 Halpern 点极限；{[}R17{]} 和 {[}R19{]}
的快速随机点极限不等于指定解选择。

\textbf{剩余问题。}
证明同一无限维随机轨道强趋于选定解，同时不牺牲声明的残差阶。应先查 G-EAG
的不精确版本以及随机 Tikhonov 文献，避免把弱到强的变化直接当作缺口。

2024--2026 代表文献为
{[}R9{]}、{[}R14{]}、{[}R17{]}。关系为\textbf{背景及可借用技巧}。它改变了用户最初的算法，不列为当前主线。

\subsection{6.
用户课题的五篇最近邻}\label{ux7528ux6237ux8bfeux9898ux7684ux4e94ux7bc7ux6700ux8fd1ux90bb}

\subsubsection{6.1
VR-FKM：最直接的有限和竞争者}\label{vr-fkmux6700ux76f4ux63a5ux7684ux6709ux9650ux548cux7adeux4e89ux8005}

{[}R16{]} 是 Tran-Dinh 单作者预印本，首稿 2024，本报告使用
\textbf{arXiv:2406.02413v4，2026-09-09}。Theorem 2 已证明
\(k^2\lVert Gx_k\rVert^2\to0\) a.s.、相邻点差的同阶小 o 和迭代点 a.s.
趋于解。SVRG / SAGA 实例的样本复杂度为
\(O(n+n^{2/3}\varepsilon^{-1})\)。

它依赖有限和及分量平均的余强制型关系，不能只凭``均值算子余强制''就套用。用户若研究一般流式
oracle，需要进一步对比 AFP。其 PDF 使用 Mathematics of Operations
Research 模板，\textbf{模板不表示正式录用}；本报告仍记预印本。

\subsubsection{6.2 AFP：一般随机 oracle 下更接近 N6
的覆盖者}\label{afpux4e00ux822cux968fux673a-oracle-ux4e0bux66f4ux63a5ux8fd1-n6-ux7684ux8986ux76d6ux8005}

{[}R17{]} 的 Theorem 2 直接处理递推误差条件，Corollary 2
包含一般无偏随机 oracle 的递增 batch
实现。零延迟、无递归误差项时，其充分条件可覆盖加权平方可和情形。因此，``从有限和推广到一般随机
oracle''不能直接作为用户课题的贡献。

它在查询点 \(y_k\) 证明快速残差，三个迭代序列有 a.s.
点极限。当前尚未完成与原始 Fast KM
的完整消元等价证明。需要比较全部系数及初始化，不能只比较图示或动量项外观。文章明示在有限维工作，并提及向
Hilbert 推广的可能性；Hilbert 本身不是充分的新颖性证据。

\subsubsection{6.3 VAPEG：Fast OGDA
方向的直接近邻}\label{vapegfast-ogda-ux65b9ux5411ux7684ux76f4ux63a5ux8fd1ux90bb}

{[}R19{]} 的首稿标题含 VFOG。\textbf{本报告采用
arXiv:2508.16791v2，2026-09-09，标题为 New Accelerated
Past-Extragradient Methods with Variance Reduction for Generalized
Equations。} 方法名为 VAPEG，a.s. 定理为 Theorem 3.2。

取集合值部分为零、单调特例后，该定理给出原始点和前一查询点的残差范数
\(o(k^{-1})\) a.s.，并有点极限。Corollary 3.4 的 \(b_k\asymp k^{3+\nu}\)
实现给出 \(O(\varepsilon^{-4-\nu})\) oracle
成本。其对数批量选项另给期望复杂度，不能自动宣称满足无限可和的 a.s.
条件。

用户需要证明原始 Fast OGDA
的噪声传播与该方法不同，且差异产生明确收益。仅强调使用 optimistic 或
past-gradient 项无法避开此文。

\subsubsection{6.4 stochastic
iKM：复杂度标杆与确定性等价线索}\label{stochastic-ikmux590dux6742ux5ea6ux6807ux6746ux4e0eux786eux5b9aux6027ux7b49ux4ef7ux7ebfux7d22}

{[}R22{]} 的 Theorem 1 给末次 RMS
残差及近最优样本复杂度，预先确定总迭代数。它不是用户所需的无限轨道 a.s.
定理。

Appendix A Proposition 1 将确定性 Fast KM 表示为 iKM
的特定参数实例。但``存在一种等价参数''与``该参数恰好满足随机主定理''是两个命题。必须检查参数及随机更新是否同时一致，才能转用结论。此文对
oracle 成本的比较价值高，不能因为它没有 a.s. 点极限就忽略其竞争性。

\subsubsection{6.5 stochastic heavy-ball
dynamics：连续时间的直接背景}\label{stochastic-heavy-ball-dynamicsux8fdeux7eedux65f6ux95f4ux7684ux76f4ux63a5ux80ccux666f}

{[}R27{]} Theorem 3.5 给 stochastic Fast OGDA 动力系统的 \(o(s^{-1})\)
a.s.
残差与弱点极限，噪声满足时间加权平方可积条件。它说明连续时间随机快速轨道已有成熟结构。

它没有自动解决离散算法。离散化产生的误差、缓存 oracle
的相关性与实际采样成本仍需另证。用户可借用能量函数的组织方式，但不能将连续时间结论替代离散证明。

\subsection{7.
建议研究入口及一个必要的反例}\label{ux5efaux8baeux7814ux7a76ux5165ux53e3ux53caux4e00ux4e2aux5fc5ux8981ux7684ux53cdux4f8b}

\subsubsection{7.1
冻结原始递推}\label{ux51bbux7ed3ux539fux59cbux9012ux63a8}

对 \(V=I-T\)，以同一个缓存 \(g_{k-1}\) 更新原始 Fast KM：

\[
x_{k+1}=x_k+\frac{k}{k+\alpha}(x_k-x_{k-1})-\frac{s\alpha}{2(k+\alpha)}g_k-\frac{sk}{k+\alpha}(g_k-g_{k-1}),\qquad g_k=Vx_k+e_k.
\]

在抽取 \(e_k\) 之前形成 \(\mathcal F_k\)。此时 \(x_k\)、\(e_{k-1}\)
已可测，新误差满足条件均值为零，旧误差不再是相对于 \(\mathcal F_k\)
的新鞅差。这一点必须直接进入 Lyapunov 漂移计算。

建议的实际顺序为：

\begin{enumerate}
\def\labelenumi{\arabic{enumi}.}
\tightlist
\item
  对 {[}R17{]} 的三序列公式消元，逐项比较状态、初值、系数和噪声。对
  {[}R16{]}、{[}R22{]} 做同样核对。得到明确的覆盖或不覆盖判断。
\item
  复现确定性 Fast KM
  的能量不等式，再保留旧误差交叉项，确定所需可和权重。先验证能量非负与条件漂移，再使用随机极限工具。
\item
  只有在得到真实差异后，才进入 5.2 或
  5.5。二维线性试验用于发现错误，不能替代 a.s. 证明。
\end{enumerate}

\subsubsection{7.2
诊断命题：二阶矩界不能支持任意降低衰减门槛}\label{ux8bcaux65adux547dux9898ux4e8cux9636ux77e9ux754cux4e0dux80fdux652fux6301ux4efbux610fux964dux4f4eux8870ux51cfux95e8ux69db}

\textbf{本报告推导，非文献新颖性主张。} 取标量
\(V(x)=x\)、\(s=1\)、\(\alpha>2\)。它满足余强制条件，也对应非扩张映射
\(T=0\)。上述递推精确化为

\[
x_{k+1}=c_kx_k-c_ke_k+d_ke_{k-1},\qquad c_k=\frac{k+\alpha/2}{k+\alpha},\quad d_k=\frac{k}{k+\alpha}.
\]

\textbf{持续方差的障碍。} 若新噪声条件方差等于 \(\sigma^2>0\)，则

\[
\operatorname{Var}(x_{k+1}\mid\mathcal F_k)=c_k^2\sigma^2.
\]

由于
\(c_k\to1\)，末次均方残差不可能趋于零。这个结论针对上述原始更新，不否定使用衰减步长、VR
或其他参数的算法。方差下界本身也不是一般 a.s. 不可能性的完整证明。

\textbf{稀有误差的 a.s. 障碍。} 固定
\(2<q\le3\)，令各步噪声独立，\(e_k\) 分别以概率 \(p_k/2\) 取 \(1/k\) 与
\(-1/k\)，其余概率取零，其中 \(p_k=k^{2-q}\)。于是

\[
\mathbb E[e_k\mid\mathcal F_k]=0,\qquad \mathbb E[e_k^2\mid\mathcal F_k]=k^{-q}.
\]

给定历史，令 \(m_k=c_kx_k+d_ke_{k-1}\)。数 \(m_k-c_k/k\) 和
\(m_k+c_k/k\) 中至少一个的绝对值不小于 \(c_k/k\)。因此

\[
\mathbb P\!\left(k\lvert x_{k+1}\rvert\ge c_k\mid\mathcal F_k\right)\ge p_k/2.
\]

\(\sum p_k=\infty\)，条件 Borel--Cantelli 引理说明该事件 a.s.
发生无穷多次。由于 \(c_k\ge1/2\)，\(k\lvert x_{k+1}\rvert\)
不能趋于零。因而，\textbf{仅有
\(\sigma_k^2=O(k^{-q})\)、\(2<q\le3\)，不足以保证这一原始递推的末次
\(o(k^{-1})\) a.s. 速率。}

这排除了初步方案中``只用二阶矩将幂次门槛从 \(q>3\) 直接放宽到
\(q>2\)''的宽泛目标。它没有证明所有算法都需要相同门槛，也没有排除轻尾或平均化样本具有额外结构时的改进。第
5.2 节因此明确加入尾部假设，并仍标为待验证。

\subsubsection{7.3
本轮方案细化的状态}\label{ux672cux8f6eux65b9ux6848ux7ec6ux5316ux7684ux72b6ux6001}

已完成问题锚定、最近邻核对、反例诊断及候选路线筛选。\texttt{research-refine}
所指定的外部 \texttt{gpt-5.5} 复核调用返回
403，提示模型未启用；\textbf{没有获得独立评分，也没有 READY
判定}。保留完整失败返回和自行修订记录。报告中的候选课题均不以该复核作为证明或背书。

\subsection{8. 文献与版本}\label{ux6587ux732eux4e0eux7248ux672c}

正文用短编号定位，以下保留英文标题。标为预印本的条目不推断录用状态。正式版与
arXiv 定理编号可能不同；正文明确版本时按该版本编号。

\textbf{{[}R1{]}} Patrick L. Combettes, Jean-Christophe Pesquet.
\emph{Stochastic Quasi-Fejér Block-Coordinate Fixed Point Iterations
with Random Sweeping}. 2015, \textbf{SIAM Journal on Optimization},
25(2):1221--1248.
\href{https://doi.org/10.1137/140971233}{DOI:10.1137/140971233}；\href{https://arxiv.org/abs/1404.7536v4}{arXiv:1404.7536v4}。正文采用该版本
Corollary 2.7。

\textbf{{[}R2{]}} Othmane Sebbouh, Robert M. Gower, Aaron Defazio.
\emph{Almost sure convergence rates for Stochastic Gradient Descent and
Stochastic Heavy Ball}. 2021, \textbf{COLT}, PMLR 134.
\href{https://arxiv.org/abs/2006.07867}{arXiv:2006.07867}；\href{https://proceedings.mlr.press/v134/sebbouh21a.html}{会议正式页}。正文采用会议版
Theorem 13、Corollary 17。

\textbf{{[}R3{]}} Jun Liu, Ye Yuan. \emph{On Almost Sure Convergence
Rates of Stochastic Gradient Methods}. 2022, \textbf{COLT}, PMLR
178:2963--2983.
\href{https://arxiv.org/abs/2202.04295}{arXiv:2202.04295}；\href{https://proceedings.mlr.press/v178/liu22d.html}{会议正式页}。

\textbf{{[}R4{]}} Julien Hermant, Marien Renaud, Jean-François Aujol,
Charles Dossal, Aude Rondepierre. \emph{Gradient correlation is a key
ingredient to accelerate SGD with momentum}. 2025, \textbf{ICLR}.
\href{https://arxiv.org/abs/2410.07870}{arXiv:2410.07870}。正文使用本地
Zotero 的 ICLR 2025 正式稿 Theorems 3--4，未把 2026 v3
的可能编号变化混入。

\textbf{{[}R5{]}} Hedy Attouch, Jalal Fadili, Vyacheslav Kungurtsev.
\emph{The Stochastic Ravine Accelerated Gradient Method with General
Extrapolation Coefficients}. 2026, \textbf{Journal of Optimization
Theory and Applications}, 209:5；2024 预印本。
\href{https://doi.org/10.1007/s10957-026-02940-4}{DOI:10.1007/s10957-026-02940-4}；\href{https://arxiv.org/abs/2403.04860}{arXiv:2403.04860}。正文按期刊版编号。

\textbf{{[}R6{]}} Saverio Salzo. \emph{The iterates of FISTA converge
even under inexact computations and stochastic gradients}. 2025
首稿，2026-07-30
v3，\textbf{预印本}。\href{https://arxiv.org/abs/2511.12665v3}{arXiv:2511.12665v3}。

\textbf{{[}R7{]}} Mario Bravo, Roberto Cominetti. \emph{Stochastic
Fixed-Point Iterations for Nonexpansive Maps: Convergence and Error
Bounds}. 2024, \textbf{SIAM Journal on Control and Optimization}.
\href{https://doi.org/10.1137/22M1515550}{DOI:10.1137/22M1515550}；\href{https://arxiv.org/abs/2208.04193v3}{arXiv:2208.04193v3}。

\textbf{{[}R8{]}} Daniel Cortild, Coralia Cartis, Juan Peypouquet.
\emph{Stochastic Krasnoselskii-Mann Iterations: Convergence without
Uniformly Bounded Variance}. 2026-09-11
v3，\textbf{预印本}。\href{https://arxiv.org/abs/2604.22581v3}{arXiv:2604.22581v3}。

\textbf{{[}R9{]}} Mario Bravo, Juan Pablo Contreras. \emph{Stochastic
Halpern iteration in normed spaces and applications to reinforcement
learning}. 2026, \textbf{Mathematical Programming}；2024 首稿。
\href{https://doi.org/10.1007/s10107-026-02345-1}{DOI:10.1007/s10107-026-02345-1}；\href{https://arxiv.org/abs/2403.12338v4}{arXiv:2403.12338v4}。本报告定理编号按
v4 正文。

\textbf{{[}R10{]}} Radu Ioan Boţ, Dang-Khoa Nguyen. \emph{Fast
Krasnosel'skiĭ--Mann Algorithm with a Convergence Rate of the Fixed
Point Iteration of \(o(1/k)\)}. 2023, \textbf{SIAM Journal on Numerical
Analysis}.
\href{https://doi.org/10.1137/22M1504305}{DOI:10.1137/22M1504305}；\href{https://arxiv.org/abs/2206.09462v2}{arXiv:2206.09462v2}。

\textbf{{[}R11{]}} Mario Bravo, Roberto Cominetti, Matías Pavez-Signé.
\emph{Rates of convergence for inexact Krasnosel'skii-Mann iterations in
Banach spaces}. 2017
预印本，正式卷期本次未复核。\href{https://arxiv.org/abs/1705.09340v2}{arXiv:1705.09340v2}。定理编号按此预印本。

\textbf{{[}R12{]}} Felix Lieder. \emph{On the convergence rate of the
Halpern-iteration}. 2021 卷期、2020 online，\textbf{Optimization
Letters}.
\href{https://doi.org/10.1007/s11590-020-01617-9}{DOI:10.1007/s11590-020-01617-9}。

\textbf{{[}R13{]}} TaeHo Yoon, Ernest K. Ryu. \emph{Accelerated
Algorithms for Smooth Convex-Concave Minimax Problems with \(O(1/k^2)\)
Rate on Squared Gradient Norm}. 2021, \textbf{ICML}, PMLR
139:12098--12109.
\href{https://arxiv.org/abs/2102.07922}{arXiv:2102.07922}；\href{https://proceedings.mlr.press/v139/yoon21d.html}{会议正式页}。

\textbf{{[}R14{]}} Radu Ioan Boţ, Enis Chenchene. \emph{Extragradient
Method with Flexible Anchoring: Strong Convergence and Fast Residual
Decay}. 2026, \textbf{SIAM Journal on Optimization}, 36(3):1420--1445.
\href{https://doi.org/10.1137/24M1703562}{DOI:10.1137/24M1703562}。

\textbf{{[}R15{]}} Radu Ioan Boţ, Enis Chenchene, Jalal Fadili.
\emph{Fast Krasnoselskii--Mann Method with Overrelaxation and
Preconditioners}. 2026, \textbf{IMA Journal of Numerical Analysis},
drag042.
\href{https://doi.org/10.1093/imanum/drag042}{DOI:10.1093/imanum/drag042}；\href{https://arxiv.org/abs/2411.18574}{arXiv:2411.18574}。正文按期刊版
Theorem 2。

\textbf{{[}R16{]}} Quoc Tran-Dinh. \emph{Variance-Reduced Fast
Krasnoselkii-Mann Methods for Finite-Sum Root-Finding Problems}. 2024
首稿，2026-09-09
v4，\textbf{预印本}。\href{https://arxiv.org/abs/2406.02413v4}{arXiv:2406.02413v4}。

\textbf{{[}R17{]}} Nghia Nguyen-Trung, Quoc Tran-Dinh. \emph{A Class of
Accelerated Fixed-Point-Based Methods with Delayed Inexact Oracles and
Its Applications}. 2025-12-15
v1，\textbf{预印本}。\href{https://arxiv.org/abs/2512.13547v1}{arXiv:2512.13547v1}。

\textbf{{[}R18{]}} Quoc Tran-Dinh. \emph{VFOSA: Variance-Reduced Fast
Operator Splitting Algorithms for Generalized Equations}. 2025,
\textbf{Journal of Machine Learning Research}, 26(254):1--68.
\href{https://www.jmlr.org/papers/v26/25-0500.html}{JMLR
正式全文及元数据}。官方页未列 DOI；对应 arXiv 编号待核实。文章存在性与
Theorem 14 已从正式全文核对，自动标题匹配仍为待核验状态。

\textbf{{[}R19{]}} Quoc Tran-Dinh, Nghia Nguyen-Trung. \emph{New
Accelerated Past-Extragradient Methods with Variance Reduction for
Generalized Equations}. 2025 首稿，2026-09-09
v2，\textbf{预印本}。\href{https://arxiv.org/abs/2508.16791v2}{arXiv:2508.16791v2}。v1
的 VFOG 标题与 v2 的 VAPEG 标题属于同一条版本记录，不计为两篇独立成果。

\textbf{{[}R20{]}} Radu Ioan Boţ, Ernö Robert Csetnek, Dang-Khoa Nguyen.
\emph{Fast Optimistic Gradient Descent Ascent (OGDA) Method in
Continuous and Discrete Time}. 2025 卷期、2023
online，\textbf{Foundations of Computational Mathematics}, 25:163--222.
\href{https://doi.org/10.1007/s10208-023-09636-5}{DOI:10.1007/s10208-023-09636-5}；\href{https://arxiv.org/abs/2203.10947}{arXiv:2203.10947}。正文按期刊版
Theorems 16--17。

\textbf{{[}R21{]}} Xufeng Cai, Chaobing Song, Cristóbal Guzmán, Jelena
Diakonikolas. \emph{Stochastic Halpern Iteration with Variance Reduction
for Stochastic Monotone Inclusions}. 2022, \textbf{NeurIPS}；2023
v4。\href{https://arxiv.org/abs/2203.09436v4}{arXiv:2203.09436v4}。正文按
v4 的 Theorem 3.1、Corollary 3.4、Theorem 4.1。

\textbf{{[}R22{]}} Tong Yang, Tao Jiang, Yuejie Chi, Ashok Cutkosky, Lin
Xiao. \emph{Stochastic Inertial Krasnosel'skii-Mann Iteration Achieves
Near-Optimal Sample Complexity}. 2026-09-22
v1，\textbf{预印本}。\href{https://arxiv.org/abs/2609.28543v1}{arXiv:2609.28543v1}。

\textbf{{[}R23{]}} TaeHo Yoon, Nicolas Loizou. \emph{Direct Acceleration
of Stochastic Root-Finding Without Variance Reduction and
Regularization}. 2026，\textbf{预印本；作者 CV 列为 NeurIPS
2026，会议官方状态待复核}。\href{https://arxiv.org/abs/2608.12043}{arXiv:2608.12043}。

\textbf{{[}R24{]}} TaeHo Yoon, Nicolas Loizou. \emph{Dual-Anchor
Acceleration Is Near-Optimal for Stochastic Monotone Root-Finding}.
2026-09-28
v1，\textbf{预印本}。\href{https://arxiv.org/abs/2609.36033v1}{arXiv:2609.36033v1}。

\textbf{{[}R25{]}} Jun-Hyun Kim, Ahmet Alacaoglu. \emph{Improving the
Last-Iterate Guarantees of Anytime Algorithms for Stochastic Monotone
Variational Inequalities}. 2026-09-24
v2，\textbf{预印本}。\href{https://arxiv.org/abs/2609.15257v2}{arXiv:2609.15257v2}。Jun-Hyun
Kim 与初始人物线索中的 Donghwan Kim 不是同一个人。

\textbf{{[}R26{]}} Jelena Diakonikolas. \emph{Solving Stochastic
Fixed-Point Equations with High Probability}.
2026，\textbf{预印本}。\href{https://arxiv.org/abs/2607.09097}{arXiv:2607.09097}。本次实际读取
arXiv HTML 的 Theorem 1 与 Corollaries 1--3。

\textbf{{[}R27{]}} Radu Ioan Boţ, Chiara Schindler. \emph{Long-time
analysis of stochastic heavy ball dynamics for convex optimization and
monotone equations}. 2026, \textbf{Discrete and Continuous Dynamical
Systems}, 51:439--474.
\href{https://doi.org/10.3934/dcds.2026035}{DOI:10.3934/dcds.2026035}。

\textbf{{[}R28{]}} Jisun Park, Ernest K. Ryu. \emph{Exact Optimal
Accelerated Complexity for Fixed-Point Iterations}. 2022, \textbf{ICML}.
\href{https://arxiv.org/abs/2201.11413v2}{arXiv:2201.11413v2}。

\textbf{{[}R29{]}} Yekini Shehu. \emph{Adaptivity, Anchoring, and the
Exact Oracle Complexity of Stochastic Fixed-Point Iterations}.
2026，\textbf{预印本}。\href{https://arxiv.org/abs/2610.00295v1}{arXiv:2610.00295v1}。正文只引用其
Theorem 4.8 所属的 Hilbert 复杂度方向，不以标题替代对所有模型的证明。

\textbf{{[}R30{]}} Jelena Diakonikolas, Cristóbal Guzmán, David
Martínez-Rubio. \emph{Oracle Complexity of Stochastic Fixed-Point
Equations with Nonexpansive Maps}.
2026，\textbf{预印本}。\href{https://arxiv.org/abs/2609.09524}{arXiv:2609.09524}。本次元数据核实，正文定理待核实，未据此填写数值速率。

\textbf{{[}R31{]}} Patrick L. Combettes, Jean-Christophe Pesquet.
\emph{Stochastic Quasi-Fejér Block-Coordinate Fixed Point Iterations
With Random Sweeping II: Mean-Square and Linear Convergence}. 2019,
\textbf{Mathematical Programming}, 174:433--451.
\href{https://arxiv.org/abs/1704.08083}{arXiv:1704.08083}。用作研究者代表作，不向本报告一般非扩张模型转用其线性率。

\textbf{{[}R32{]}} Juan Pablo Contreras, Roberto Cominetti.
\emph{Optimal error bounds for non-expansive fixed-point iterations in
normed spaces}. 2023, \textbf{Mathematical Programming}, 199:343--374.
\href{https://doi.org/10.1007/s10107-022-01830-7}{DOI:10.1007/s10107-022-01830-7}；\href{https://arxiv.org/abs/2108.10969}{arXiv:2108.10969}。

\textbf{{[}R33{]}} Saverio Salzo, Silvia Villa. \emph{Inexact and
Accelerated Proximal Point Algorithms}. 2012, \textbf{Journal of Convex
Analysis}, 19(4):1167--1192.
\href{https://doi.org/10.68381/jca19062}{DOI:10.68381/jca19062}；\href{https://journalofconvexanalysis.com/articles/jca19062/}{期刊正式页}。

\textbf{{[}R34{]}} Ahmet Alacaoglu, Yura Malitsky. \emph{Stochastic
Variance Reduction for Variational Inequality Methods}. 2022,
\textbf{COLT}, PMLR 178:778--816.
\href{https://arxiv.org/abs/2102.08352}{arXiv:2102.08352}；\href{https://proceedings.mlr.press/v178/alacaoglu22a.html}{会议正式页}。

\textbf{{[}R35{]}} Radu Ioan Boţ, Ernö Robert Csetnek, Christopher
Hendrich. \emph{Inertial Douglas--Rachford splitting for monotone
inclusion problems}. 2015, \textbf{Applied Mathematics and Computation},
256:472--487.
\href{https://doi.org/10.1016/j.amc.2015.01.017}{DOI:10.1016/j.amc.2015.01.017}；\href{https://arxiv.org/abs/1403.3330}{arXiv:1403.3330}。

\subsection{9.
核查范围与使用限制}\label{ux6838ux67e5ux8303ux56f4ux4e0eux4f7fux7528ux9650ux5236}

本轮按三个层次交叉核对：先用本地讨论笔记和 Zotero
确认研究对象，再阅读原文定理和参数条件，最后用 arXiv / DOI /
官方发表页面复核作者、版本与发表信息。这是三个核查环节，\textbf{不表示每个结论都有三个相互独立的证明来源}。

原论文定理是速率的依据；作者主页用来确认身份与任职；元数据只确认书目，不能证明定理正确。核心算法的递推等价、弱点极限推广和候选更弱噪声定理仍为待验证事项。报告没有对预印本证明做逐行形式审查。

原始线索中的 \texttt{2609.09152} 不对应 stochastic
iKM；本次采用已经核实的 {[}R22{]}
\texttt{2609.28543}。此前抓取到的一个同姓 Yoon 的 ICML PDF 实为
autoencoding 论文，已从证据链排除，EAG 改用已核对作者和标题的
\texttt{yoon21d} 正式全文。

配套研究目录保存候选文献、自动核验状态、主页访问记录、核心论文知识库、覆盖关系、反例诊断和方案细化日志。报告正文可独立阅读；这些辅助文件用于后续研究追踪。外部复核的服务失败不影响已保存的本地文献证据，但独立评审仍未完成。

\end{document}
